\documentclass[11pt,a4paper]{article}
\usepackage[spanish,es-nodecimaldot]{babel}
\usepackage[utf8]{inputenc}
\usepackage[T1]{fontenc}
\usepackage{lmodern}
\usepackage{microtype}
\usepackage{amsmath,amssymb,amsthm,bm}
\usepackage{geometry}
\usepackage{booktabs,tabularx,array,longtable}
\usepackage{graphicx,float,subcaption}
\usepackage{xcolor}
\usepackage{hyperref}
\usepackage{listings}
\usepackage{enumitem}
\usepackage{fancyhdr}
\geometry{margin=2.35cm,headheight=15pt}
\hypersetup{colorlinks=true,linkcolor=blue!50!black,urlcolor=blue!50!black,citecolor=blue!50!black}
\pagestyle{fancy}
\fancyhf{}
\fancyhead[L]{Reto III.3: asistencia gravitacional}
\fancyhead[R]{Sol--Júpiter--Voyager 2}
\fancyfoot[C]{\thepage}
\setlength{\parskip}{.45em}
\setlength{\parindent}{0pt}
\setlist[itemize]{topsep=3pt,itemsep=2pt}
\definecolor{codebg}{RGB}{244,247,250}
\definecolor{codeblue}{RGB}{20,65,112}
\lstdefinestyle{python}{language=Python,basicstyle=\ttfamily\footnotesize,
  keywordstyle=\color{codeblue}\bfseries,commentstyle=\color{teal!60!black},
  stringstyle=\color{orange!65!black},breaklines=true,
  backgroundcolor=\color{codebg},frame=single,showstringspaces=false,
  numbers=left,numberstyle=\tiny\color{gray},xleftmargin=1.2em,
  literate={á}{{\'a}}1 {é}{{\'e}}1 {í}{{\'i}}1 {ó}{{\'o}}1 {ú}{{\'u}}1 {ñ}{{\~n}}1 {ü}{{\"u}}1}
\lstset{style=python}
\newcommand{\vect}[1]{\bm{#1}}
\newcommand{\dif}{\mathop{}\!\mathrm{d}}
\newcommand{\SOI}{R_{\mathrm{SOI}}}
\newcommand{\muS}{\mu_\odot}
\newcommand{\muJ}{\mu_{J}}
\newcommand{\rJ}{\vect r_{J}}
\newcommand{\vJ}{\vect v_{J}}
\newcommand{\rV}{\vect r}
\newcommand{\vV}{\vect v}
\newcommand{\fign}[3]{\begin{figure}[htbp]\centering\IfFileExists{figuras/#1}{\includegraphics[width=#2\textwidth]{figuras/#1}}{\fbox{\parbox{.85\textwidth}{Archivo opcional: \texttt{figuras/#1}}}}\caption{#3}\end{figure}}
\begin{document}
\begin{titlepage}
\centering
\vspace*{2cm}
{\Large Mecánica clásica y física computacional\par}
\vspace{1.2cm}
{\Huge\bfseries Reto III.3\par}
\vspace{.5cm}
{\LARGE Diseño y validación de una asistencia gravitacional\par}
\vspace{.65cm}
{\Large Caso histórico: Voyager 2 en Júpiter (1979)\par}
\vspace{1.1cm}
\rule{.8\textwidth}{.6pt}\par
\vspace{.5cm}
{\large\bfseries Sistema dinámico: Sol, Júpiter y sonda\par}
{\normalsize Dos aproximaciones: cónicas empalmadas e integración gravitacional\par}
\vspace{1.4cm}
\begin{minipage}{.83\textwidth}
\textbf{Alcance:} se prescriben las efemérides heliocéntricas de Júpiter y se considera la sonda de masa despreciable. El radio $9.58\,\mathrm{UA}$ sirve solo como referencia geométrica de salida hacia la región de Saturno; \emph{no se incluyen efemérides ni aceleración gravitatoria de Saturno}.
\end{minipage}
\vfill
{\large Documento metodológico y resultados reproducibles\par}
{\large Noviembre de 2026\par}
\end{titlepage}
\tableofcontents
\newpage

\section{Resumen, problema y objetivos}
El propósito es reconstruir el sobrevuelo de la sonda Voyager 2 por Júpiter mediante las efemérides de JPL Horizons, modelar la desviación gravitatoria y demostrar cómo cambia su trayectoria en el sistema heliocéntrico. Se comparan dos descripciones que comparten exactamente la misma condición inicial de entrada: (i) una integración de dos cuerpos alrededor de Júpiter limitada por una esfera de influencia, unida a arcos heliocéntricos keplerianos, y (ii) una propagación restringida con la acción simultánea del Sol y Júpiter.

\textbf{Objetivo físico:} determinar qué energía y qué momento angular se conservan, en qué referencia y bajo qué hipótesis. \textbf{Objetivo numérico:} cuantificar el error de la aproximación de cónicas empalmadas y su dependencia del radio de cambio de modelo. \textbf{Objetivo geométrico:} establecer si, después del sobrevuelo, la sonda cruza el radio heliocéntrico $R_{\mathrm{obj}}=9.58\,\mathrm{UA}$. Esto último es una condición necesaria de alcance radial, \emph{no una predicción de encuentro con Saturno}.

\section{Marco teórico}
\subsection{Sistema de referencia, notación y unidades}
Trabajamos en coordenadas heliocéntricas eclípticas J2000 del archivo Horizons, con posiciones en $\mathrm{km}$, velocidades en $\mathrm{km\,s^{-1}}$, tiempos en segundos y época astronómica TDB. Denotamos con $\rV(t)$ y $\vV(t)$ el estado de Voyager y con $\rJ(t)$ y $\vJ(t)$ el de Júpiter, ambos respecto del Sol. La masa de la sonda es despreciable frente a las masas de los primarios.

Definimos el estado relativo jovicéntrico por
\begin{equation}\label{eq:rel}
\vect\rho=\rV-\rJ,\qquad \vect w=\vV-\vJ,
\qquad \rho=\lVert\vect\rho\rVert.
\end{equation}
La ecuación~\eqref{eq:rel} es una traslación galileana instantánea de posiciones y velocidades. El sistema centrado en Júpiter, al acelerarse con su movimiento orbital, \emph{no es inercial} durante un intervalo largo. Por tanto no debe aplicarse directamente la ecuación jovicéntrica pura cuando se quieren incluir todas las fuerzas.

En el código se emplean:
\begin{align*}
\muS &= 1.32712440018\times 10^{11}\ \mathrm{km^3\,s^{-2}},\\
\muJ &= 1.26686534\times 10^{8}\ \mathrm{km^3\,s^{-2}},\\
1\ \mathrm{UA}&=149\,597\,870.7\ \mathrm{km},\qquad 1\ \mathrm{día}=86\,400\ \mathrm{s},\\
a_J&=778.57\times10^6\ \mathrm{km},\\
M_\odot &= 1.98847\times10^{30}\ \mathrm{kg},\qquad M_J=1.89813\times10^{27}\ \mathrm{kg}.
\end{align*}
Se utilizan valores constantes para los parámetros gravitacionales y un semieje mayor representativo de Júpiter. Las efemérides no se sustituyen por órbitas circulares al calcular la trayectoria nominal.

\subsection{Problema de dos cuerpos en un campo central}
Para la posición relativa a un cuerpo de parámetro gravitacional $\mu=GM$, el modelo kepleriano satisface
\begin{equation}\label{eq:kepler}
\ddot{\vect r}=-\frac{\mu}{r^3}\vect r,\qquad
\vect F=-\frac{GMm}{r^3}\vect r.
\end{equation}
La fuerza es central: siempre paralela a $\vect r$. La energía específica (energía por unidad de masa de la sonda) es
\begin{equation}\label{eq:energy}
\varepsilon=\frac{\dot{\vect r}\cdot\dot{\vect r}}2-\frac{\mu}{r}.
\end{equation}
Derivando respecto al tiempo,
\begin{align}
\frac{\dif\varepsilon}{\dif t}
 &=\dot{\vect r}\cdot\ddot{\vect r}+\mu\frac{\vect r\cdot\dot{\vect r}}{r^3}\\
 &=\dot{\vect r}\cdot\left(-\frac\mu{r^3}\vect r\right)
 +\mu\frac{\vect r\cdot\dot{\vect r}}{r^3}=0.
\end{align}
En consecuencia, $\varepsilon$ es constante en un problema de dos cuerpos autónomo. El momento angular específico es el \emph{vector}
\begin{equation}\label{eq:h}
\vect h=\vect r\times\dot{\vect r}.
\end{equation}
Su derivada es
\begin{equation}
\dot{\vect h}=\dot{\vect r}\times\dot{\vect r}+
\vect r\times\ddot{\vect r}=\vect 0,
\end{equation}
porque $\ddot{\vect r}$ es paralelo a $\vect r$. La conservación \emph{vectorial} fija el plano orbital, no solamente el módulo $h=\lVert\vect h\rVert$.

\subsection{Energía, excentricidad y elementos osculatrices}
El vector de excentricidad es
\begin{equation}
\vect e=\frac{\vect v\times\vect h}{\mu}-\frac{\vect r}{r},
\qquad e=\lVert\vect e\rVert.
\end{equation}
La integración del problema central conduce a la ecuación polar de una cónica
\begin{equation}\label{eq:conic}
r(\nu)=\frac{p}{1+e\cos\nu},\qquad p=\frac{h^2}{\mu}.
\end{equation}
La energía determina el semieje mayor
\begin{equation}
\varepsilon=-\frac{\mu}{2a},\qquad a=-\frac{\mu}{2\varepsilon}.
\end{equation}
Para $\varepsilon<0$ la órbita es elíptica ($a>0$), con afelio $r_a=a(1+e)$. Si $\varepsilon>0$, la cónica es hiperbólica ($a<0$, $e>1$) y \emph{no hay afelio finito}. La inclinación, en eclíptica J2000, se calcula mediante
\begin{equation}
i=\arccos\left(\frac{h_z}{\lVert\vect h\rVert}\right).
\end{equation}
Los elementos calculados durante el encuentro se llaman \emph{osculatrices}: representan la cónica instantánea asociada al estado, no una trayectoria exacta si hay perturbaciones.

\subsection{Dispersión hiperbólica por gravitación}
En el sobrevuelo jovicéntrico ideal, la energía relativa es positiva y puede definirse la rapidez asintótica por
\begin{equation}
\varepsilon_J=\frac12w^2-\frac{\muJ}{\rho}>0,
\qquad v_\infty=\sqrt{2\varepsilon_J}.
\end{equation}
Del momento angular $h_J=\rho\,w_\perp$, muy lejos de Júpiter se obtiene el parámetro de impacto
\begin{equation}\label{eq:b}
b=\frac{h_J}{v_\infty}.
\end{equation}
A partir de la energía y del momento angular,
\begin{equation}\label{eq:ecc}
e=\sqrt{1+\frac{2\varepsilon_J h_J^2}{\muJ^2}},\qquad
r_p=\frac{h_J^2}{\muJ(1+e)}.
\end{equation}
Las asíntotas de la hipérbola satisfacen $\cos\nu_\infty=-1/e$. El cambio de dirección entre las velocidades asintóticas es
\begin{equation}\label{eq:gamma}
\gamma=2\arcsin\left(\frac1e\right)=2\arctan\left(\frac{\muJ}{bv_\infty^2}\right).
\end{equation}
Según Berg~\cite{berg2023}, ecuaciones (26)--(27), se define el parámetro adimensional
\begin{equation}\label{eq:D}
D=\frac{\muJ}{b v_\infty^2},\qquad \gamma=2\arctan D.
\end{equation}
Así, $D\ll1$ representa una dispersión débil y $D\gg1$ una gran desviación. Es importante separar $\gamma$, \emph{ángulo entre asíntotas}, del giro medido entre velocidades sobre la \emph{frontera finita} de la SOI: estas cantidades no tienen por qué coincidir.

\subsection{Por qué una sonda gana energía heliocéntrica}
Sean $\vect w_-$ y $\vect w_+$ las velocidades relativas a Júpiter antes y después del sobrevuelo, tomadas suficientemente lejos del planeta, y sea $\vJ$ su velocidad heliocéntrica aproximadamente constante durante la maniobra corta. En esa aproximación,
\begin{equation}
\vect v_-=\vJ+\vect w_-,\qquad
\vect v_+=\vJ+\vect w_+,\qquad |\vect w_-|=|\vect w_+|=v_\infty.
\end{equation}
Por tanto, la diferencia de energía cinética heliocéntrica específica es
\begin{align}\label{eq:transfer}
\Delta K_\odot
&=\frac12|\vect v_+|^2-\frac12|\vect v_-|^2\\
&=\vJ\cdot(\vect w_+-\vect w_-)
+\frac12\bigl(|\vect w_+|^2-|\vect w_-|^2\bigr)\\
&=\boxed{\vJ\cdot(\vect w_+-\vect w_-)}.
\end{align}
La sonda puede ganar energía heliocéntrica porque Júpiter está en movimiento; en un tratamiento de masas finitas se intercambian cantidades muy pequeñas de energía y momento con el planeta. Esto \textbf{no viola} la conservación de la energía del sistema físico completo.

Berg~\cite{berg2023} introduce la geometría plana con la rapidez planetaria $u_p$, la rapidez asintótica relativa $u_i$, el ángulo jovicéntrico de entrada $\beta$ y el giro $\gamma$. La ecuación (32) del artículo equivale a
\begin{multline}\label{eq:berg32}
\Delta u=\sqrt{u_p^2+u_i^2-2u_pu_i\cos(\beta+\gamma)}\\
-\sqrt{u_p^2+u_i^2-2u_pu_i\cos\beta}.
\end{multline}
Usando
\begin{equation}
\cos\gamma=\frac{1-D^2}{1+D^2},\qquad
\sin\gamma=\frac{2D}{1+D^2},
\end{equation}
se transforma la ecuación~\eqref{eq:berg32} en la ecuación (37) de Berg,
\begin{multline}\label{eq:berg37}
\Delta u=\Biggl[u_p^2+u_i^2-2u_pu_i\left(
\frac{1-D^2}{1+D^2}\cos\beta-\frac{2D}{1+D^2}\sin\beta\right)\Biggr]^{1/2}\\
-\sqrt{u_p^2+u_i^2-2u_pu_i\cos\beta}.
\end{multline}
Para relacionar la entrada heliocéntrica con la jovicéntrica (ecuaciones (39)--(42)), el plano se orienta con eje tangencial al movimiento de Júpiter y eje radial:
\begin{align}
u'_{m,x}&=u'_m\sin\alpha,& u_{m,x}&=u'_{m,x}-u_p,\\
u_i&=\sqrt{u_{m,x}^2+(u'_m\cos\alpha)^2},&
\cos\beta&=\frac{|u_{m,x}|}{u_i}.
\end{align}
\textbf{Advertencia de validación:} estas expresiones suponían una orientación concreta del giro, un plano de dispersión y velocidad planetaria aproximadamente tangencial. En el código se comprueban además las velocidades asintóticas \emph{vectoriales 3D} y la ganancia entre superficies finitas de la SOI. No se debe obligar a que estas cuatro evaluaciones sean exactamente iguales.

\subsection{Esfera de influencia: Laplace y Hill}
La aproximación de cónicas empalmadas necesita una frontera convencional. Se adopta el radio de Laplace,
\begin{equation}\label{eq:soi}
\SOI=a_J\left(\frac{M_J}{M_\odot}\right)^{2/5},
\end{equation}
obtenido en el tratamiento clásico de la dominancia de perturbaciones relativas. Como escala alternativa se calcula el radio de Hill para una órbita aproximadamente circular,
\begin{equation}\label{eq:hill}
R_{\mathrm{Hill}}=a_J\left(\frac{M_J}{3M_\odot}\right)^{1/3}.
\end{equation}
Con las constantes declaradas,
\begin{equation}
\SOI=48\,219\,247.59\ \mathrm{km},\qquad
R_{\mathrm{Hill}}=53\,152\,808.54\ \mathrm{km}
\quad (R_{\mathrm{Hill}}/\SOI\approx1.1023).
\end{equation}
Ninguno de los radios representa una superficie física abrupta donde desaparezca exactamente la influencia solar. Por ello se realiza un estudio de sensibilidad para distintos radios de empalme.

\section{Formulación del Modelo 1: cónicas empalmadas}
\subsection{Estado de entrada y cambio de marco}
La distancia real de la sonda a Júpiter se define por
\begin{equation}
d_J(t)=\bigl\lVert\rV^{\rm Horizons}(t)-\rJ^{\rm Horizons}(t)\bigr\rVert.
\end{equation}
El instante de entrada $t_{\rm in}$ se obtiene encontrando el primer cruce descendente $d_J(t)=\SOI$. Se inicializan las variables jovicéntricas en el cruce:
\begin{equation}\label{eq:entry}
\vect\rho(t_{\rm in})=\rV(t_{\rm in})-\rJ(t_{\rm in}),
\quad \vect w(t_{\rm in})=\vV(t_{\rm in})-\vJ(t_{\rm in}).
\end{equation}
Para el segmento interior se integra \emph{únicamente}
\begin{equation}\label{eq:model1}
\dot{\vect\rho}=\vect w,
\qquad
\dot{\vect w}=-\frac{\muJ}{\rho^3}\vect\rho.
\end{equation}
El evento terminal de salida corresponde al primer cruce ascendente
\begin{equation}
g_{\rm out}(t)=\lVert\vect\rho(t)\rVert-\SOI=0,
\quad \frac{\dif g_{\rm out}}{\dif t}>0.
\end{equation}
En el instante calculado $t_{\rm out}^{(1)}$ se transforman nuevamente las condiciones:
\begin{equation}\label{eq:exit}
\rV(t_{\rm out}^{(1)})=\rJ(t_{\rm out}^{(1)})+\vect\rho(t_{\rm out}^{(1)}),
\quad
\vV(t_{\rm out}^{(1)})=\vJ(t_{\rm out}^{(1)})+\vect w(t_{\rm out}^{(1)}).
\end{equation}
Se integran segmentos solares independientes antes y después del encuentro:
\begin{equation}
\ddot{\rV}=-\frac{\muS}{r^3}\rV.
\end{equation}
La trayectoria queda definida por tres soluciones empalmadas en $t_{\rm in}$ y $t_{\rm out}^{(1)}$. El estado de entrada es \emph{idéntico} para ambos modelos, mientras que sus fechas de salida pueden diferir y deben conservarse para interpretar el método.

\subsection{Conservaciones de la región jovicéntrica}
En el intervalo interior del Modelo 1, de la ecuación~\eqref{eq:model1} se deduce directamente
\begin{equation}\label{eq:inv1}
\boxed{\frac{\dif\varepsilon_J}{\dif t}=0,\qquad
\frac{\dif\vect h_J}{\dif t}=\vect0},
\quad
\varepsilon_J=\frac12 w^2-\frac{\muJ}{\rho},\quad
\vect h_J=\vect\rho\times\vect w.
\end{equation}
No se exige que la energía de Kepler alrededor del Sol sea constante \emph{durante} el segmento jovicéntrico, porque se está cambiando entre modelos y referencias. Tampoco se exige la continuidad de la aceleración en la frontera: sí son continuas posiciones y velocidades por construcción.

\section{Formulación del Modelo 2: Sol, Júpiter y sonda}
\subsection{Ecuación heliocéntrica restringida con término indirecto}
El Modelo 2 no integra la evolución gravitatoria de todos los cuerpos desde condiciones iniciales propias. La trayectoria real de Júpiter, $\rJ(t)$, es una función prescrita por Horizons. Se propaga únicamente una sonda de masa despreciable. Debido a que el origen está ligado al Sol, debe introducirse la aceleración de ese origen por Júpiter, llamada \emph{término indirecto}:
\begin{equation}\label{eq:model2}
\boxed{\ddot{\rV}=
-\muS\frac{\rV}{r^3}
-\muJ\left[
\frac{\rV-\rJ(t)}{|\rV-\rJ(t)|^3}
+\frac{\rJ(t)}{|\rJ(t)|^3}
\right].}
\end{equation}
El primer término corresponde a la atracción solar, el segundo a la atracción \emph{directa} de Júpiter, y el tercero a la corrección indirecta por la aceleración del Sol. La sonda recibe ambas contribuciones en todo el intervalo; aquí \textbf{no existe una frontera de cambio de ecuación}.

Si Júpiter fuese un primario dinámico del problema ideal de dos cuerpos Sol--Júpiter, la segunda ecuación sería
\begin{equation}
\ddot{\rJ}=-(\muS+\muJ)\frac{\rJ}{|\rJ|^3}.
\end{equation}
Pero en este proyecto $\rJ(t)$ se extrae de un sistema solar real y puede contener perturbaciones omitidas del modelo. La formulación es, por tanto, una \emph{aproximación restringida forzada por efemérides}, no una simulación autónoma completa de tres cuerpos.

\subsection{Fuerza diferencial, marea solar y derivadas de invariantes}
De $\vect\rho=\rV-\rJ$ se sigue
\begin{align}\label{eq:rel_acc}
\ddot{\vect\rho}
&=-\muJ\frac{\vect\rho}{\rho^3}+\vect a_{\rm marea}+\vect\delta_a,\\
\vect a_{\rm marea}
&= -\muS\left(\frac{\rV}{r^3}-\frac{\rJ}{|\rJ|^3}\right),\\
\vect\delta_a
&=-\muS\frac{\rJ}{|\rJ|^3}
-\muJ\frac{\rJ}{|\rJ|^3}
-\ddot{\rJ}_{\rm Horizons}.
\end{align}
La corrección $\vect\delta_a$ cuantifica la diferencia entre la aceleración joviana coherente con el modelo ideal Sol--Júpiter y la aceleración de las efemérides prescritas. En una órbita planetaria ideal compatible, $\vect\delta_a=\vect0$.

Para el Modelo 2, las variaciones de las cantidades jovicéntricas se obtienen sin aproximaciones adicionales a partir de la ecuación relativa:
\begin{align}\label{eq:derivE}
\boxed{\frac{\dif\varepsilon_J}{\dif t}=\vect w\cdot
(\vect a_{\rm marea}+\vect\delta_a)},\\
\boxed{\frac{\dif\vect h_J}{\dif t}=\vect\rho\times
(\vect a_{\rm marea}+\vect\delta_a)}.
\end{align}
El cociente adimensional que se traza es el diagnóstico de marea \emph{solar} respecto a la gravedad joviana directa,
\begin{equation}\label{eq:tidal}
\eta(t)=\frac{|\vect a_{\rm marea}|}{\muJ/\rho^2}.
\end{equation}
Este cociente no incluye $\vect\delta_a$ y no es un criterio universal de validez por sí solo; además, las pequeñas perturbaciones integradas pueden producir grandes diferencias a largo plazo.

Por su parte, la energía específica de Kepler heliocéntrica y el momento angular heliocéntrico satisfacen
\begin{align}
\varepsilon_\odot&=\frac12v^2-\frac{\muS}{r},
&\frac{\dif\varepsilon_\odot}{\dif t}
&=\vV\cdot\vect a_{J,\rm total},\\
\vect h_\odot&=\rV\times\vV,
&\frac{\dif\vect h_\odot}{\dif t}
&=\rV\times\vect a_{J,\rm total},
\end{align}
con
\begin{equation}
\vect a_{J,\rm total}=-\muJ\left(
\frac{\rV-\rJ}{|\rV-\rJ|^3}+\frac{\rJ}{|\rJ|^3}\right).
\end{equation}
Por tanto $\varepsilon_\odot$ y $\vect h_\odot$ tampoco son invariantes exactos del estado de la sonda bajo la fuerza combinada. No debe confundirse esta afirmación con la conservación total que tendría un \emph{sistema físico autónomo} de masas finitas.

\subsection{Constante de Jacobi y alcance de la comprobación}
En el \emph{problema restringido circular} de tres cuerpos, usando un sistema sinódico uniformemente rotante, existen un potencial efectivo $U_{\rm ef}$ y una constante
\begin{equation}\label{eq:jacobi}
C_J=2U_{\rm ef}(\vect r_{\rm rot})-\lVert\vect v_{\rm rot}\rVert^2.
\end{equation}
Para un sistema baricéntrico circular, con $\vect r_B,\vect v_B$ y velocidad angular constante $n\widehat{\vect n}$, esto se puede escribir en las variables inerciales baricéntricas como
\begin{equation}\label{eq:cstar}
C_J=2\left(\frac{\muS}{r_\odot}+\frac{\muJ}{r_J}\right)
-\lVert\vect v_B\rVert^2+2n\,\widehat{\vect n}\cdot(\vect r_B\times\vect v_B),
\end{equation}
con las distancias a los dos primarios tomadas desde sus posiciones baricéntricas. En el código $\vect r_B=\vect r-\chi\rJ$, $\vect v_B=\vect v-\chi\vJ$ y $\chi=\muJ/(\muS+\muJ)$, usando un eje y velocidad angular de referencia construidos localmente. Como Horizons no representa exactamente una órbita circular uniforme de Sol--Júpiter, el valor calculado se llama $C^*(t)$, \emph{cuasi-Jacobi}, y se usa exclusivamente como diagnóstico. Una prueba separada con órbita circular consistente permite comprobar que la fórmula tiende a conservarse en el caso ideal.

\section{Procedimiento computacional completo}
\subsection{Lectura y sincronización de JPL Horizons}
Los únicos archivos necesarios son \texttt{horizons\_jupiter.csv} y \texttt{horizons\_voyager-2.csv}. La rutina \texttt{horizons\_csv} acepta tanto el formato tabular limpio como la exportación con marcadores \texttt{\$\$SOE} y \texttt{\$\$EOE}. Comprueba el encabezado cuando está disponible: \emph{centro Sol}, \emph{unidades KM-S} y \emph{eclíptica J2000}. Se rechazan fechas inconsistentes.

Para interpolar \emph{posición y velocidad de forma compatible}, se utiliza un spline cúbico de Hermite, no dos interpoladores independientes. Para cada componente de posición entre dos muestras $t_k,t_{k+1}$, con $s=(t-t_k)/\Delta t$, se interpola
\begin{multline}\label{eq:hermite}
r(t)=(2s^3-3s^2+1)r_k+ (s^3-2s^2+s)\Delta t\,v_k\\
+(-2s^3+3s^2)r_{k+1}+(s^3-s^2)\Delta t\,v_{k+1},
\end{multline}
con $v(t)=\dif r/\dif t$. Este método preserva los valores y derivadas tabulados pero no inventa precisión física superior a la cadencia original de Horizons.

\subsection{Eventos geométricos: entrada, periapsis, salida y radio objetivo}
Primero se detectan índices donde $d_J-\SOI$ cambia de positivo a negativo para la entrada y viceversa para la salida. Cada cruce se refina mediante \texttt{brentq}; la fecha de máxima aproximación se calcula minimizando $d_J(t)$ con \texttt{minimize\_scalar} en el intervalo de la SOI. Para la proyección se busca el primer cruce ascendente de
\begin{equation}
g_{\rm obj}(t)=\lVert\rV(t)\rVert-R_{\rm obj},\quad
R_{\rm obj}=9.58\times149\,597\,870.7\ \mathrm{km}.
\end{equation}
Este evento no comprueba fase angular ni posición de Saturno, pues ese planeta no forma parte del modelo.

\subsection{Integración ODE y tolerancias}
En ambos modelos se escribe un sistema de primer orden de seis variables $\vect y=(x,y,z,v_x,v_y,v_z)$ y se usa \texttt{scipy.integrate.solve\_ivp} con \texttt{method='DOP853'} (Runge--Kutta explícito de orden alto), \texttt{dense\_output=True}, $\mathrm{rtol}=3\times10^{-10}$ y tolerancias absolutas separadas para posición $10^{-3}\,\mathrm{km}$ y velocidad $2\times10^{-11}\,\mathrm{km/s}$. El paso máximo nominal es de $0.75$ días; el integrador acorta automáticamente los pasos cerca del periapsis.

\begin{lstlisting}[caption={Ecuaciones de propagación: código abreviado, idéntico a las fuerzas usadas.}]
def solar_rhs(phys):
    def f(t, y):
        r, v = y[:3], y[3:]
        return np.r_[v, -phys.mu_sun*r / np.linalg.norm(r)**3]
    return f

def jovian_rhs(phys):
    def f(t, y):
        rho, w = y[:3], y[3:]
        return np.r_[w, -phys.mu_jupiter*rho / np.linalg.norm(rho)**3]
    return f

def planet_rhs(phys, eph):
    def f(t, y):
        r, v = y[:3], y[3:]
        rj = eph["jupiter"].state(t).r
        rho = r-rj
        a = -phys.mu_sun*r/np.linalg.norm(r)**3
        a -= phys.mu_jupiter*(rho/np.linalg.norm(rho)**3
                               + rj/np.linalg.norm(rj)**3)
        return np.r_[v, a]
    return f
\end{lstlisting}
\textbf{Nota:} el listado muestra las ecuaciones físicas de manera legible; las funciones ejecutables, la gestión de eventos y el programa \emph{completo} están incorporados en el cuaderno y disponibles en \texttt{asistencia\_gravitacional\_sol\_jupiter.py}.

\subsection{Esquema de flujo y relación con funciones Python}
\begin{enumerate}
\item \texttt{synchronize\_data}: leer únicamente Júpiter y Voyager y verificar mismo intervalo y malla temporal de Horizons.
\item \texttt{encounter\_real}: localizar $t_{\rm in}$, $t_{\rm peri}$ y $t_{\rm out}$ a partir de $d_J(t)$ y $\SOI$.
\item \texttt{conics\_orbit}: anclar el mismo estado de entrada y construir los tres arcos del Modelo 1: solar anterior, jovicéntrico, solar posterior.
\item \texttt{forced\_orbit}: con la misma entrada, integrar la ecuación~\eqref{eq:model2} hacia atrás y hacia adelante sin empalmes físicos.
\item \texttt{hyperbola\_geometry}: calcular $\varepsilon_J,h_J,v_\infty,b,e,D,\gamma,r_p$ a partir del estado jovicéntrico de entrada.
\item \texttt{berg\_parameters}: trasladar asintotas relativas a base tangencial--radial, calcular $\alpha$, $\beta$ y las ecuaciones (32), (37) y la comparación 3D.
\item \texttt{two\_body\_invariants}, \texttt{tidal\_ratio}, \texttt{jacobi\_proxy}: generar diagnósticos de conservación y perturbación.
\item \texttt{sensitivity\_study}: recalcular la entrada para distintos radios SOI y repetir ambas propagaciones con la \emph{misma} condición inicial en cada experimento.
\item \texttt{first\_radius\_time}: encontrar el primer cruce de $9.58\,\mathrm{UA}$ en cada modelo.
\item \texttt{make\_plots}: exportar once figuras de Matplotlib; \texttt{plotly\_animation\_2d} exporta el visor 2D; \texttt{generar\_panel3d} exporta las vistas interactivas jovicéntrica y heliocéntrica por Canvas 2D, sin WebGL.
\item \texttt{main}: escribir archivos CSV, resumen JSON y resultados, y mostrar verificaciones. El cuaderno Colab agrupa los productos en un ZIP descargable \emph{desde Colab}.
\end{enumerate}

\subsection{Definición de error y comparación en el mismo instante}
Para dos estados calculados en una misma época $t$ se utiliza
\begin{equation}
\delta_r(t)=\lVert\rV_1(t)-\rV_2(t)\rVert,
\qquad \delta_v(t)=\lVert\vV_1(t)-\vV_2(t)\rVert.
\end{equation}
Frente a Horizons se define de modo análogo $\delta_r^{\rm real}$ y $\delta_v^{\rm real}$. Si el Modelo 1 sale de la SOI antes que el Modelo 2, sus estados se evalúan \emph{ambos a la fecha de salida del Modelo 1}; no se deben restar dos posiciones tomadas en tiempos diferentes. El error de proyección se evalúa también en la misma \emph{fecha real de cruce} del radio $9.58$ UA. Las diferencias M1--M2 son el efecto de cambiar las ecuaciones aproximadas; las diferencias frente a Horizons también reflejan fuerzas que los modelos simplificados omiten.

\subsection{Sensibilidad al radio de empalme}
Se repite el procedimiento para
\begin{equation}
R_{\rm empalme}=f\,\SOI,
\quad f\in\left\{0.50,0.65,0.80,1.00,\frac{R_{\rm Hill}}{\SOI}\right\}.
\end{equation}
En cada caso se determina una nueva fecha de entrada sobre Horizons y se inicializan ambos modelos con \emph{ese mismo estado real}. La comparación no equivale a controlar todos los errores físicos, pero permite aislar la sensibilidad del esquema de cónicas empalmadas ante una decisión convencional del analista.

\section{Resultados obtenidos con los archivos del proyecto}
\subsection{Reconstrucción del encuentro real}
Los valores obtenidos con los CSV de Horizons y la interpolación Hermite son:
\begin{center}\begin{tabular}{ll}\toprule
Magnitud & Resultado\\\midrule
Entrada SOI Laplace & 6 de mayo de 1979, 11:34:09 TDB\\
Periapsis interpolado & 9 de julio de 1979, 22:26:54 TDB\\
Distancia mínima interpolada & $720\,894$ km\\
Salida SOI real & 12 de septiembre de 1979, 20:40:14 TDB\\
SOI Laplace / Hill & $48.219$ / $53.153$ millones de km\\
Cruce real de $9.58$ UA & 23 de agosto de 1981, 07:21:45 TDB\\
\bottomrule\end{tabular}\end{center}
El valor de periapsis depende de la interpolación y de la cadencia de los registros. No se presenta como una reconstrucción de navegación de alta precisión de la misión.

\subsection{Elementos heliocéntricos: antes y después}
Los elementos de la sonda reconstruidos a partir de Horizons en las fronteras reales de Laplace son:
\begin{center}\begin{tabular}{lrr}\toprule
Magnitud & Antes de la SOI & Después de la SOI\\\midrule
Energía específica $\varepsilon_\odot$ ($\mathrm{km^2/s^2}$) & $-121.446$ & $30.091$\\
Semieje mayor osculatriz $a$ (UA) & $3.652$ & $-14.741$\\
Excentricidad $e$ & $0.7371$ & $1.3406$\\
Afelio (UA) & $6.345$ & no existe\\
Inclinación eclíptica $i$ (grados) & $5.1205$ & $2.5799$\\\bottomrule\end{tabular}\end{center}
En la entrada el afelio osculatriz solar es inferior a $9.58\,\mathrm{UA}$; al salir, la cónica solar osculatriz es abierta ($\varepsilon_\odot>0$). Esto muestra la posibilidad de viajar a radios mayores sin impulso propulsivo adicional, siempre que se tenga en cuenta el movimiento de Júpiter.

\subsection{Parámetros analíticos de la hipérbola}
En el estado de entrada de Laplace:
\begin{center}\begin{tabular}{ll}\toprule
Parámetro & Valor\\\midrule
$\varepsilon_J$ & $29.43869\ \mathrm{km^2/s^2}$\\
$h_J$ & $1.269885\times10^7\ \mathrm{km^2/s}$\\
$v_\infty$ & $7.67316\ \mathrm{km/s}$\\
$b$ & $1\,654\,970\ \mathrm{km}$\\
$D$ & $1.300145$\\
Excentricidad $e_J$ & $1.26158$\\
Giro asintótico $\gamma$ & $104.869^\circ$\\
Periapsis de la cónica osculatriz & $562\,842\ \mathrm{km}$\\
$\alpha$ y $\beta$ proyectados según Berg & $53.337^\circ$ y $51.499^\circ$\\
Ganancia Berg (32) y (37) & $10.0076\ \mathrm{km/s}$\\
Ganancia asintótica vectorial 3D & $9.7525\ \mathrm{km/s}$\\\bottomrule\end{tabular}\end{center}
El periapsis de la hipérbola osculatriz pura difiere del periapsis de las efemérides reales: ello ilustra el efecto de fuerzas no incluidas en la aproximación. El ángulo de $104.869^\circ$ corresponde a la \emph{cónica jovicéntrica osculatriz en la SOI}, no al ángulo histórico nominal de unos $80^\circ$ mencionado por Berg como aproximación geométrica; no deben identificarse sin revisar las definiciones y condiciones iniciales.

\subsection{Ejemplo completo de sustitución numérica}
Para mostrar de dónde surgen los valores y permitir reproducirlos manualmente, tomamos el estado relativo \emph{interpolado} en $t_{\rm in}$ y las constantes de Júpiter. El código calcula primero
\begin{align*}
\varepsilon_J &=\frac12\lVert\vect w_0\rVert^2-\frac{\muJ}{\lVert\vect\rho_0\rVert}
 \approx29.43869\ \mathrm{km^2/s^2},\\
h_J&=\lVert\vect\rho_0\times\vect w_0\rVert
 \approx1.269885\times10^7\ \mathrm{km^2/s}.
\end{align*}
Una vez obtenidos estos \emph{dos invariantes osculatrices} se suceden las sustituciones
\begin{align*}
v_\infty&=\sqrt{2(29.43869)}\approx7.67316\ \mathrm{km/s},\\
b&=\frac{1.269885\times10^7}{7.67316}
 \approx1.65497\times10^6\ \mathrm{km},\\
D&=\frac{1.26686534\times10^8}{(1.65497\times10^6)(7.67316)^2}
 \approx1.300145,\\
\gamma&=2\arctan(1.300145)\approx104.869^\circ,\\
e_J&=\sqrt{1+\frac{2(29.43869)(1.269885\times10^7)^2}
 {(1.26686534\times10^8)^2}}\approx1.26158,\\
r_{p,\mathrm{osc}}&=\frac{(1.269885\times10^7)^2}{
(1.26686534\times10^8)(1+1.26158)}\approx5.62842\times10^5\ \mathrm{km}.
\end{align*}
El mismo procedimiento aplicado a cada estado heliocéntrico de Horizons da
\begin{align*}
 a_{\rm antes}&=-\frac{\muS}{2\varepsilon_{\odot,\rm antes}}
 =-\frac{1.32712440018\times10^{11}}{2(-121.446)}
 \approx 5.464\times10^8\ \mathrm{km}\approx3.652\ \mathrm{UA},\\
 r_{a,\rm antes}&=a_{\rm antes}(1+e_{\rm antes})
 \approx3.652(1+0.7371)=6.345\ \mathrm{UA},\\
 a_{\rm despues}&=-\frac{1.32712440018\times10^{11}}{2(30.0909)}
 \approx-14.741\ \mathrm{UA},\\
 \Delta\varepsilon_\odot&=\varepsilon_{\odot,\rm despues}-
 \varepsilon_{\odot,\rm antes}
 =30.0909-(-121.4460)\approx151.537\ \mathrm{km^2/s^2}.
\end{align*}
Estos cálculos son \emph{instantáneos u osculatrices}; las posiciones, velocidades, cruces y errores a fechas posteriores requieren la propagación numérica. Las cifras redondeadas aquí no deben reemplazar los valores de doble precisión usados por el código.

\subsection{Cuantificación de los dos modelos}
Los dos modelos empiezan en la misma época y estado de entrada. Para cada uno se mide la energía entre esa entrada y su \emph{propia} salida, y el error respecto a Horizons se calcula en esa misma fecha de salida:
\begin{center}\begin{tabular}{lrr}\toprule
Resultado & Modelo 1 & Modelo 2\\\midrule
$\Delta\varepsilon_\odot$ ($\mathrm{km^2/s^2}$) & $160.4714$ & $152.2056$\\
Error de posición al salir (km) & $6\,845\,903$ & $469\,583$\\
Error de velocidad al salir ($\mathrm{km/s}$) & $1.1256$ & $0.08715$\\
Cruce radial de $9.58$ UA & 10--08--1981 & 19--08--1981\\\bottomrule\end{tabular}\end{center}
Para comparación, Horizons da $\Delta\varepsilon_\odot\approx151.537\,\mathrm{km^2/s^2}$ entre \emph{sus propias} fronteras. La coincidencia de energía no es una prueba suficiente de exactitud orbital: se deben verificar también vector posición y vector velocidad.

\subsection{¿Qué se conserva realmente?}
En el Modelo 1, calculando la diferencia extrema de energía jovicéntrica dentro del tramo de dos cuerpos y las variaciones relativas del momento angular vectorial,
\begin{equation}
\operatorname{rango}(\varepsilon_J)\approx6.57\times10^{-8}\ \mathrm{km^2/s^2},
\quad \max_t\frac{|\vect h_J(t)-\vect h_J(0)|}{|\vect h_J(0)|}
\approx1.15\times10^{-10}.
\end{equation}
Son residuos de integración numérica. En el Modelo 2 la variación total de $\varepsilon_J$ en la ventana común llega a $0.275\,\mathrm{km^2/s^2}$ y el cambio relativo vectorial máximo de $\vect h_J$ es cercano a $15.2\%$. Estas variaciones son físicas dentro del modelo forzado y reflejan la presencia de aceleraciones no centrales jovicéntricas. No indican que se pierda energía en el universo.

El cociente $\eta$ de la ecuación~\eqref{eq:tidal} varía en la ventana analizada entre $7.76\times10^{-7}$ y $0.4267$. La perturbación relativa es muy pequeña cerca del planeta, pero considerable en la periferia de la SOI. Los umbrales ilustrativos $\delta_r\geq10^3\,\mathrm{km}$, $\delta_r\geq10^5\,\mathrm{km}$ y $\delta_v\geq0.1\,\mathrm{km/s}$ aparecen aproximadamente el 9 de mayo, 12 de junio y 4 de julio de 1979, respectivamente. \emph{Son convenciones diagnósticas, no un instante universal de ruptura de equivalencia.}

\subsection{Resultados del estudio de sensibilidad}
\begin{center}\begin{tabular}{rrr}\toprule
Radio $f\,\SOI$ & Error M1--M2 a salida (millones km) & Diferencia proyectada a 9.58 UA (millones km)\\\midrule
$0.50$ & $0.374$ & $39.358$\\
$0.65$ & $1.182$ & $32.296$\\
$0.80$ & $2.771$ & $40.280$\\
$1.00$ & $7.190$ & $76.182$\\
Hill ($1.102$) & $10.598$ & $102.087$\\\bottomrule\end{tabular}\end{center}
El error proyectado mínimo de este conjunto discreto se obtiene para $f=0.65$; no hay garantía de que sea un óptimo físico universal, ni necesariamente el mejor valor para otro objetivo, época o sistema.

\subsection{Figuras del cálculo}\label{sec:figuras}
\fign{01_trayectoria_jovicentrica.png}{.96}{Reconstrucción de la trayectoria cerca de Júpiter para Horizons, cónicas empalmadas y propagación simultánea.}
\fign{04_invariantes_energia_momento.png}{.96}{Energías específicas y momentos angulares en los dos marcos de referencia.}
\fign{06_marea_y_error.png}{.90}{Relación entre marea solar, aceleración joviana y divergencia de los modelos.}
\fign{08_trayectoria_radio_objetivo.png}{.85}{Proyección heliocéntrica; $9.58$ UA es un radio de referencia, no una órbita planetaria calculada.}
\fign{10_sensibilidad_soi.png}{.93}{Errores frente a los distintos radios de empalme, comparados en tiempos comunes.}

\section{Descripción y organización del código}
El proyecto se puede ejecutar por completo en Google Colab sin modificar archivos de sistema ni recurrir a GPU. El cuaderno incluye el motor científico en una celda propia, el análisis, la generación de once figuras PNG, las tablas CSV, la animación Plotly 2D y dos HTML independientes con Canvas 2D. La cámara interactiva se centra de forma predeterminada en Júpiter; se puede seleccionar seguimiento de sonda o Sol, y trasladar la vista con \emph{Mayús + arrastre} o botón derecho. No se utiliza \texttt{go.Scatter3d}, que requeriría WebGL.

\textbf{Arquitectura:} \texttt{Physics} agrupa parámetros físicos, \texttt{Settings} tolerancias y opciones, \texttt{Ephemeris} encapsula datos e interpolación, \texttt{State} describe cada estado físico, \texttt{Encounter} agrupa eventos geométricos, y \texttt{Orbit} evalúa tramos integrados. Funciones separadas implementan fuerzas, integraciones, métricas, diagnósticos y figuras. Los resultados se exportan en \texttt{comparacion\_modelos.csv}, \texttt{elementos\_osculatrices.csv}, \texttt{sensibilidad\_radio\_soi.csv}, \texttt{diagnosticos\_temporales.csv}, \texttt{resumen\_fisico.json} y \texttt{telemetria\_dashboard\_*.csv}.

\textbf{Procedimiento reproducible en Colab:}
\begin{enumerate}
\item Abrir \texttt{Voyager2\_RetoIII3\_SOLO\_SOL\_JUPITER.ipynb}.
\item Ejecutar la celda de dependencias y cargar únicamente los CSV de Júpiter y Voyager 2.
\item Ejecutar la celda del motor físico, y después la celda de \texttt{main()}.
\item Ejecutar las dos celdas de visualización HTML, una jovicéntrica y otra heliocéntrica.
\item Ejecutar la última celda para generar un ZIP y \emph{descargarlo desde Colab}. El archivo LaTeX y su PDF de lectura pueden conservarse en la misma carpeta del proyecto y se incluyen como documentación externa.
\end{enumerate}

\subsection{Limitaciones y comprobaciones de consistencia}
\begin{itemize}
\item Las efemérides históricas de Horizons incluyen fuerzas y características dinámicas omitidas en los dos modelos; el error contra Horizons \emph{no mide exclusivamente} el error de cónicas empalmadas. Para comparar aproximaciones en igualdad de condiciones, también se computa M1--M2.
\item El movimiento de Júpiter se toma de Horizons, mientras el modelo de la sonda solo incorpora gravedad solar y joviana. Esto es una aproximación forzada, no una integración simultánea autoconsistente de todas las masas.
\item Se trata el Sol como centro heliocéntrico con corrección indirecta; esta corrección evita mezclar sin aviso un origen acelerado con ecuaciones inerciales.
\item El parámetro de SOI y la malla temporal son decisiones de modelado. Las fechas subdiarias calculadas por interpolación no deben confundirse con mediciones originales diarias.
\item La conservación de Jacobi requiere una órbita circular coherente y un sistema rotante uniforme. Con efemérides reales solo se grafica $C^*(t)$ aproximada; una autoprueba con órbita circular ideal recupera deriva relativa del orden de $10^{-10}$.
\item Cruzar el radio $9.58\,\mathrm{UA}$ no prueba una cita orbital con Saturno. No se usan ni su posición ni su gravedad en el proyecto final.
\end{itemize}

\section{Conclusiones}
\begin{enumerate}[label=\textbf{\arabic*.}]
\item La dispersión por fuerza central proporciona una comprensión analítica del encuentro: la dirección de la velocidad relativa cambia con un giro que depende de $D=\muJ/(bv_\infty^2)$; su energía relativa y el momento angular vectorial se conservan en la descripción jovicéntrica de dos cuerpos.
\item La asistencia gravitacional no genera energía de la nada. El cambio de rapidez heliocéntrica proviene de la transformación de velocidades y del intercambio de energía orbital con Júpiter: $\Delta K_\odot\simeq\vJ\cdot(\vect w_+-\vect w_-)$.
\item El Modelo 1 reproduce el mecanismo esencial con un costo numérico bajo y conserva $\varepsilon_J$ y $\vect h_J$ en el segmento de la SOI. Sin embargo, su hipótesis elimina la marea solar durante un intervalo largo, lo que se acumula como error de trayectoria.
\item El Modelo 2, limitado a Sol--Júpiter--sonda, reproduce mejor las efemérides reales de salida: aproximadamente $0.470$ millones de km de error posicional frente a $6.846$ millones de km para el Modelo 1. La discrepancia no se explica por el integrador, sino principalmente por el modelo físico empleado.
\item No existe un instante único y universal cuando las aproximaciones ``dejan de ser equivalentes''. La respuesta depende de la tolerancia de error aceptable, del radio de empalme, de la duración de la predicción y de la magnitud comparada.
\item La energía heliocéntrica osculatriz se vuelve positiva después del paso joviano y la trayectoria alcanza el radio geométrico $9.58\,\mathrm{UA}$. Este resultado satisface la condición radial pedida sin introducir interacciones gravitatorias adicionales.
\end{enumerate}

\section{Referencias}
\begin{thebibliography}{9}
\bibitem{berg2023} P. Berg, ``The fundamental concepts of the gravity-assist manoeuvre'', \emph{European Journal of Physics}, vol.~44, art.~025002, 2023. DOI: \href{https://doi.org/10.1088/1361-6404/acaad8}{10.1088/1361-6404/acaad8}. Se utilizó la traducción al español suministrada para el proyecto, especialmente secciones 2--5 y ecuaciones (26)--(27), (32), (37), (39)--(42).
\bibitem{horizons} NASA/JPL, \emph{Horizons On-Line Ephemeris System}. Archivos CSV históricos de Júpiter y Voyager 2 aportados para este proyecto; marco de referencia y unidades según sus encabezados.
\bibitem{scipy} SciPy, documentación de \texttt{solve\_ivp} (método DOP853), \texttt{CubicHermiteSpline} y \texttt{brentq}; implementaciones utilizadas por el programa.
\end{thebibliography}
\end{document}
